% GENERATED FILE. Edit canonical lessons, then run npm run generate:books.
% canonical-source-hash: fnv1a64:ee278f9f4ad09070
% canonical-lessons: JA-C472-zairyou, JA-C472-choumiryou, JA-C472-negi, JA-C472-komugiko, JA-C472-daizu

\chapter{Ingredients}
\label{ch:ja-ingredients}
\hlchaptermodality{472}

\begin{chapteropening}
\textbf{By the end of this chapter:} \emph{I can name ingredients and seasonings.}

Complete the last lesson of chapter 472: I can name ingredients and seasonings.
\end{chapteropening}

\section[\texorpdfstring{\ja{ざいりょう} (zairyō)}{ざいりょう (zairyō)}]{\ja{ざいりょう} (zairyō) — ingredients; materials}

\label{lesson:JA-C472-zairyou}



\begin{quote}
Before the new one: say the Japanese for to sing in a choir, then the Japanese for to compose music.
\end{quote}

\subsection*{You'll want to know: \ja{ざいりょう}}
\textbf{\ja{ざいりょう}} — \emph{zairyō} — \textquotedblleft{}ingredients; materials\textquotedblright{}.

\begin{quote}
\textbf{\ja{ばんごはんの} \ja{ざいりょうを} \ja{かいました。}} — \emph{bangohan no zairyō o kaimashita} — \textquotedblleft{}I bought the ingredients for dinner.\textquotedblright{}
\end{quote}

\begin{cousinweb}
\textbf{\ja{ざいりょう}} (\emph{zairyō}) is Sino-Japanese (kango): two Chinese characters read with their on'yomi, readings borrowed from Chinese. They are \emph{zai}, \textquotedblleft{}ingredients\textquotedblright{} (Middle Chinese \emph{dzoj}) and \emph{ryō}, \textquotedblleft{}materials\textquotedblright{} (\emph{lew}) (KANJIDIC2; Middle Chinese in Baxter's transcription).
\end{cousinweb}

\begin{grammarlens}[title={What you've built}]
One more word for describing what you lived and felt.
\end{grammarlens}

\subsection*{Your turn}
\begin{itemize}
\raggedright
  \item \emph{Say it:} \emph{zairyō}
  \item \emph{Say it:} \emph{zairyō}, once more
  \item \emph{Say it:} \emph{bangohan no zairyō o kaimashita} — I bought the ingredients for dinner.
  \item \emph{Recall:} say the Japanese for to sing in a choir, then the Japanese for to compose music, then say \emph{zairyō} again
\end{itemize}

\begin{tcolorbox}[breakable,colback=teal!4,colframe=teal!35!black,title={Before you move on}]
What does \ja{ざいりょう} mean? (\textquotedblleft{}Ingredients; materials\textquotedblright{}.) Say it once more.
\end{tcolorbox}
\section[\texorpdfstring{\ja{ちょうみりょう} (chōmiryō)}{ちょうみりょう (chōmiryō)}]{\ja{ちょうみりょう} (chōmiryō) — a seasoning}

\label{lesson:JA-C472-choumiryou}



\begin{quote}
Before the new one: say the Japanese for to compose music, then the Japanese for ingredients.
\end{quote}

\subsection*{You'll want to know: \ja{ちょうみりょう}}
\textbf{\ja{ちょうみりょう}} — \emph{chōmiryō} — \textquotedblleft{}a seasoning\textquotedblright{}.

\begin{quote}
\textbf{\ja{しおと} \ja{しょうゆは} \ja{ちょうみりょうです。}} — \emph{shio to shōyu wa chōmiryō desu} — \textquotedblleft{}Salt and soy sauce are seasonings.\textquotedblright{}
\end{quote}

\begin{cousinweb}
\textbf{\ja{ちょうみりょう}} (\emph{chōmiryō}) is Sino-Japanese (kango): three Chinese characters read with their on'yomi, readings borrowed from Chinese. They are \emph{chō}, \textquotedblleft{}prepare\textquotedblright{} (Middle Chinese \emph{dew}), \emph{mi}, \textquotedblleft{}taste\textquotedblright{} (\emph{mj+jH}), and \emph{ryō}, \textquotedblleft{}materials\textquotedblright{} (\emph{lew}) (KANJIDIC2; Middle Chinese in Baxter's transcription).
\end{cousinweb}

\begin{grammarlens}[title={What you've built}]
One more word for describing what you lived and felt.
\end{grammarlens}

\subsection*{Your turn}
\begin{itemize}
\raggedright
  \item \emph{Say it:} \emph{chōmiryō}
  \item \emph{Say it:} \emph{chōmiryō}, once more
  \item \emph{Say it:} \emph{shio to shōyu wa chōmiryō desu} — Salt and soy sauce are seasonings.
  \item \emph{Recall:} say the Japanese for to compose music, then the Japanese for ingredients, then say \emph{chōmiryō} again
\end{itemize}

\begin{tcolorbox}[breakable,colback=teal!4,colframe=teal!35!black,title={Before you move on}]
What does \ja{ちょうみりょう} mean? (\textquotedblleft{}A seasoning\textquotedblright{}.) Say it once more.
\end{tcolorbox}
\section[\texorpdfstring{\ja{ねぎ} (negi)}{ねぎ (negi)}]{\ja{ねぎ} (negi) — a green onion}

\label{lesson:JA-C472-negi}



\begin{quote}
Before the new one: say the Japanese for ingredients, then the Japanese for a seasoning.
\end{quote}

\subsection*{You'll want to know: \ja{ねぎ}}
\textbf{\ja{ねぎ}} — \emph{negi} — \textquotedblleft{}a green onion\textquotedblright{}.

\begin{quote}
\textbf{\ja{みそしるに} \ja{ねぎを} \ja{いれました。}} — \emph{misoshiru ni negi o iremashita} — \textquotedblleft{}I put green onion in the miso soup.\textquotedblright{}
\end{quote}

\begin{cousinweb}
\textbf{\ja{ねぎ}} (\emph{negi}) is a native Japanese word (wago), inherited rather than borrowed from Chinese; its kanji spelling was attached to it later (JMdict).
\end{cousinweb}

\begin{grammarlens}[title={What you've built}]
One more word for describing what you lived and felt.
\end{grammarlens}

\subsection*{Your turn}
\begin{itemize}
\raggedright
  \item \emph{Say it:} \emph{negi}
  \item \emph{Say it:} \emph{negi}, once more
  \item \emph{Say it:} \emph{misoshiru ni negi o iremashita} — I put green onion in the miso soup.
  \item \emph{Recall:} say the Japanese for ingredients, then the Japanese for a seasoning, then say \emph{negi} again
\end{itemize}

\begin{tcolorbox}[breakable,colback=teal!4,colframe=teal!35!black,title={Before you move on}]
What does \ja{ねぎ} mean? (\textquotedblleft{}A green onion\textquotedblright{}.) Say it once more.
\end{tcolorbox}
\section[\texorpdfstring{\ja{こむぎこ} (komugiko)}{こむぎこ (komugiko)}]{\ja{こむぎこ} (komugiko) — wheat flour}

\label{lesson:JA-C472-komugiko}



\begin{quote}
Before the new one: say the Japanese for a seasoning, then the Japanese for a green onion.
\end{quote}

\subsection*{You'll want to know: \ja{こむぎこ}}
\textbf{\ja{こむぎこ}} — \emph{komugiko} — \textquotedblleft{}wheat flour\textquotedblright{}.

\begin{quote}
\textbf{\ja{こむぎこと} \ja{たまごを} \ja{まぜます。}} — \emph{komugiko to tamago o mazemasu} — \textquotedblleft{}I mix the flour and the eggs.\textquotedblright{}
\end{quote}

\begin{cousinweb}
\textbf{\ja{こむぎこ}} (\emph{komugiko}) is a native Japanese compound (wago), inherited rather than borrowed from Chinese: \emph{ko}, \textquotedblleft{}small\textquotedblright{}, \emph{mugi}, \textquotedblleft{}wheat\textquotedblright{}, and \emph{ko}, \textquotedblleft{}powder\textquotedblright{} (JMdict, KANJIDIC2).
\end{cousinweb}

\begin{grammarlens}[title={What you've built}]
One more word for describing what you lived and felt.
\end{grammarlens}

\subsection*{Your turn}
\begin{itemize}
\raggedright
  \item \emph{Say it:} \emph{komugiko}
  \item \emph{Say it:} \emph{komugiko}, once more
  \item \emph{Say it:} \emph{komugiko to tamago o mazemasu} — I mix the flour and the eggs.
  \item \emph{Recall:} say the Japanese for a seasoning, then the Japanese for a green onion, then say \emph{komugiko} again
\end{itemize}

\begin{tcolorbox}[breakable,colback=teal!4,colframe=teal!35!black,title={Before you move on}]
What does \ja{こむぎこ} mean? (\textquotedblleft{}Wheat flour\textquotedblright{}.) Say it once more.
\end{tcolorbox}
\section[\texorpdfstring{\ja{だいず} (daizu)}{だいず (daizu)}]{\ja{だいず} (daizu) — soybeans}

\label{lesson:JA-C472-daizu}



\begin{quote}
Before the new one: say the Japanese for a green onion, then the Japanese for wheat flour.
\end{quote}

\subsection*{You'll want to know: \ja{だいず}}
\textbf{\ja{だいず}} — \emph{daizu} — \textquotedblleft{}soybeans\textquotedblright{}.

\begin{quote}
\textbf{\ja{なっとうは} \ja{だいずの} \ja{たべものです。}} — \emph{nattō wa daizu no tabemono desu} — \textquotedblleft{}Natto is a food made from soybeans.\textquotedblright{}
\end{quote}

\begin{cousinweb}
\textbf{\ja{だいず}} (\emph{daizu}) is Sino-Japanese (kango): two Chinese characters read with their on'yomi, readings borrowed from Chinese. They are \emph{dai}, \textquotedblleft{}large\textquotedblright{} (Middle Chinese \emph{dajH}) and \emph{zu}, \textquotedblleft{}beans\textquotedblright{} (\emph{duwH}) (KANJIDIC2; Middle Chinese in Baxter's transcription).
\end{cousinweb}

\begin{grammarlens}[title={What you've built}]
One more word for describing what you lived and felt.
\end{grammarlens}

\subsection*{Your turn}
\begin{itemize}
\raggedright
  \item \emph{Say it:} \emph{daizu}
  \item \emph{Say it:} \emph{daizu}, once more
  \item \emph{Say it:} \emph{nattō wa daizu no tabemono desu} — Natto is a food made from soybeans.
  \item \emph{Recall:} say the Japanese for a green onion, then the Japanese for wheat flour, then say \emph{daizu} again
\end{itemize}

\begin{tcolorbox}[breakable,colback=teal!4,colframe=teal!35!black,title={Before you move on}]
What does \ja{だいず} mean? (\textquotedblleft{}Soybeans\textquotedblright{}.) Say it once more.
\end{tcolorbox}
